% ---------------------------------------------------------------------------
% Key Idea paragraph for the final report.
% Drop this as a one-paragraph subsection between Problem Formulation (§II)
% and Main Results (§III), either as "\subsection{Key Idea}" or as an unnumbered
% paragraph opening §III. Patterned after Example_FinalReprot2.pdf, whose single
% well-framed key-idea paragraph carried a report with no simulations.
%
% References expected: \cite{Fiedler1973,Boyd2006,GhoshBoyd2006,Schulman2017}.
% ---------------------------------------------------------------------------

\subsection*{Key Idea}
We frame algebraic-connectivity maximization under bounded edge weights
and a communication budget as a constrained Markov decision process whose
optimum is upper-bounded by a convex semidefinite program, namely the
Fastest Distributed Linear Averaging (FDLA) relaxation of Boyd and
co-authors \cite{XiaoBoyd2004,Boyd2006,GhoshBoyd2006}. This bound
reframes an otherwise heuristic benchmarking exercise as a quantitative
suboptimality claim: for any feasible reweighting policy $\pi$ acting on
a fixed support, $\lambda_2(L(w_\pi)) \le s^\star$, where $s^\star$ is
the SDP optimum. We then show that a Proximal Policy Optimization
\cite{Schulman2017} agent observing only graph features and consensus
statistics learns a policy that attains
$\lambda_2(L(w_\pi))/s^\star \approx 0.87$, averaged over five
Erd\H{o}s--R\'enyi seeds with $n = 20$ and a tight budget $B = 25$.
Because the FDLA SDP ceases to apply when the underlying edge set is
itself a decision variable (discrete rewiring) or when the graph support
is induced by agent positions (random geometric graphs), we further
evaluate the same policy class in those two settings and report its
generalization to regimes where no convex upper bound is available.
